% ============================================================
% RESULTADOS — TFG
% Índice Sintético de Turismo Sostenible por CCAA (2016–2024)
% ============================================================

\chapter{Resultados}

\section{Características del panel procesado}

Tras la aplicación del pipeline completo (imputación, tratamiento de outliers, normalización), el panel activo contiene 52 indicadores distribuidos en tres dimensiones. La imputación MICE redujo los valores ausentes de 3.065 (32,75\%) a 118 (1,26\%). El análisis de asimetría y curtosis identificó 75 valores atípicos (índices z superiores a 3 en valor absoluto), de los cuales 5 fueron sometidos a winsorización. Los indicadores más afectados fueron IAMB0002 (porcentaje de medios de transporte menos contaminante, $|z|_{\max} = 2.68$) e ISOC0007 (índice de percepción de seguridad, $|z|_{\max} = 5.31$).

\section{Pesos obtenidos mediante PCA}

\subsection{Análisis de Nivel 1: pesos de indicadores dentro de dimensiones}

El análisis de componentes principales realizado sobre el panel apilado (171 observaciones, 52 indicadores) para derivar los pesos de Nivel 1 reveló:

\subsubsection{Dimensión Económica (ECO)}

La varianza explicada por PC1 en ECO alcanzó el 48,2\%. Los indicadores con mayor peso (loading absoluto normalizado) fueron:

\begin{table}[H]
    \centering
    \small
    \begin{tabular}{llc}
        \toprule
        \textbf{Código} & \textbf{Nombre} & \textbf{Peso} \\
        \midrule
        IECO0018 & Gasto medio por turista interno, día & 0.0473 \\
        IECO0034 & RevPAR (ingresos/habitación disponible) & 0.0462 \\
        IECO0026 & Índice de Precios Hoteleros & 0.0451 \\
        IECO0019 & ADR (tarifa media diaria) & 0.0445 \\
        IECO0027 & Variación anual IPH & 0.0438 \\
        \bottomrule
    \end{tabular}
    \caption{Top 5 indicadores económicos por peso PCA}
    \label{tab:pesos_eco_top5}
\end{table}

Los indicadores de estancia media (IECO0005, IECO0009) recibieron pesos bajos (0.009--0.010), lo que refleja que su varianza entre CCAA es reducida comparada con indicadores de tarificación y gasto.

\subsubsection{Dimensión Ambiental (ENV)}

La varianza explicada por PC1 fue 39,8\%. Los indicadores se concentraron en dos componentes dual:

\begin{table}[H]
    \centering
    \small
    \begin{tabular}{llc}
        \toprule
        \textbf{Código} & \textbf{Nombre} & \textbf{Peso} \\
        \midrule
        IAMB0006 & Porcentaje de llegadas por avión & 0.3705 \\
        IAMB0007 & Porcentaje de llegadas por mar y tierra & 0.3462 \\
        IAMB0002 & Porcentaje medios transporte menos contaminante & 0.2176 \\
        IAMB0001 & Gasto en energía less contaminante & 0.0657 \\
        \bottomrule
    \end{tabular}
    \caption{Pesos de indicadores ambientales}
    \label{tab:pesos_env}
\end{table}

Los dos primeros indicadores suman el 72\% del peso de la dimensión. Son complementarios (siempre suman 100\%), indicativo de que capturan un único eje latente: la composición modal del transporte turístico. IAMB0001 y IAMB0002 capturan energía y transporte limpio respectivamente.

\subsubsection{Dimensión Social (SOC)}

La varianza explicada por PC1 alcanzó 41,5\%. Tres indicadores de satisfacción turística concentraron el peso:

\begin{table}[H]
    \centering
    \small
    \begin{tabular}{llc}
        \toprule
        \textbf{Código} & \textbf{Nombre} & \textbf{Peso} \\
        \midrule
        ISOC0013 & Índice satisfacción alojamientos hoteleros & 0.2003 \\
        ISOC0010 & Índice satisfacción hoteles 5 estrellas & 0.1842 \\
        ISOC0005 & Índice Percepción Turística Global & 0.1784 \\
        ISOC0012 & Índice satisfacción productos turísticos & 0.0948 \\
        ISOC0011 & Ocupación laboral turística & 0.0812 \\
        \bottomrule
    \end{tabular}
    \caption{Top 5 indicadores sociales por peso PCA}
    \label{tab:pesos_soc_top5}
\end{table}

Los indicadores de percepción (ISOC0006, ISOC0007, ISOC0008) recibieron pesos menores (0.043--0.062), aunque sus distribuciones presentaban mayor asimetría.

\subsection{Análisis de Nivel 2: pesos de las dimensiones}

Los pesos derivados de la segunda iteración PCA, sobre los scores dimensionales agregados, fueron:

\begin{table}[H]
    \centering
    \begin{tabular}{lc}
        \toprule
        \textbf{Dimensión} & \textbf{Peso de Nivel 2} \\
        \midrule
        Ambiental (ENV) & 0.4517 \\
        Económica (ECO) & 0.4438 \\
        Social (SOC) & 0.1045 \\
        \bottomrule
    \end{tabular}
    \caption{Pesos de las dimensiones en el índice final}
    \label{tab:pesos_dim}
\end{table}

La dimensión social recibe un peso significativamente inferior (10,45\%) comparado con ENV y ECO (que se reparten el 89,55\% de forma casi equitativa: 45,17\% y 44,38\%). Este resultado refleja que los indicadores de percepción y satisfacción turística presentan menor varianza compartida con el eje económico-ambiental que captura PC1. 

En particular, el indicador IAMB0006 y IAMB0007, siendo duales, generan una correlación negativa que PC1 captura como un único factor dominante. El peso reducido de SOC es consistente: en años con mayor variabilidad como 2022, el peso de SOC asciende a 0.104, mientras que en años con menores contrastes (ej. 2020, año de pandemia con comportamiento homogéneo), desciende a 0.042.

\section{Estructura de correlaciones e indicadores redundantes}

El análisis de conglomerados para el año 2022 (año con máxima disponibilidad de datos) mediante clustering jerárquico (distancia euclídea, método Ward D2, $k=4$) identificó:

\begin{itemize}
    \item \textbf{Cluster 1} (30 indicadores): precios hoteleros por categoría y tiempo de antelación. Altamente cohesivo ($r > 0.80$ entre pares), reflejo de que miden esencialmente un único fenómeno: nivel tarifario del alojamiento.
    \item \textbf{Cluster 2} (8 indicadores): mezcla de gasto, estancia, transporte ambiental e índices de precios. Estructura hetereogénea, interpretable como «mercado e impacto ambiental general».
    \item \textbf{Cluster 3} (11 indicadores): variaciones de tarifa y satisfacción turística. PC1 de este cluster captura tensión entre dinamismo de precios y valoración del turista.
    \item \textbf{Cluster 4} (3 indicadores): percepciones de seguridad y clima, más un indicador económico residual. Aislado por sus distribuciones leptocúrticas extremas (Kurtosis 7.19 y 4.97).
\end{itemize}

La elevada redundancia intra-ECO (Cluster 1) justifica el empleo de PCA: sin el ajuste por capacidad discriminante, el índice habría estado dominado por variaciones de tarificación hotelera.

\section{Varianza explicada en el sistema completo}

El PCA exploratorio sobre el panel normalizado completo (171 observaciones, 52 indicadores) mostró:

\begin{table}[H]
    \centering
    \small
    \begin{tabular}{cc}
        \toprule
        \textbf{Número de PC} & \textbf{Varianza acumulada (\%)} \\
        \midrule
        PC1 & 12.52 \\
        PC1--2 & 21.85 \\
        PC1--5 & 35.42 \\
        PC1--10 & 51.63 \\
        PC1--20 & 75.48 \\
        \bottomrule
    \end{tabular}
    \caption{Varianza explicada acumulada del panel normalizado}
    \label{tab:var_acumulada}
\end{table}

El hecho de que PC1 explique solo el 12,52\% indica que la estructura de variabilidad en sostenibilidad turística es compleja y multidimensional, sin una dirección dominante única. Esto valida el enfoque PCA en dos niveles: no sería apropiado utilizarun único componente global para ponderar todos los indicadores.

\section{Análisis de robustez: comparación de cuatro casos metodológicos}

Se comparan cuatro combinaciones (Casos 1--4) sobre el período completo 2016--2024. Para facilitar la comparación se presenta el año 2024 como referencia.

\subsection{Caso 1: Min-Max + PCA + Media Aritmética}

\begin{table}[H]
    \centering
    \small
    \begin{tabular}{lrr}
        \toprule
        \textbf{Posición} & \textbf{CCAA} & \textbf{Score} \\
        \midrule
        1 & Navarra (ES-NC) & 57.68 \\
        2 & Asturias (ES-AS) & 54.48 \\
        3 & Galicia (ES-GA) & 53.64 \\
        4 & Castilla y León (ES-CL) & 52.31 \\
        5 & Cantabria (ES-CB) & 51.59 \\
        \midrule
        15 & Baleares (ES-IB) & 46.94 \\
        19 & Murcia (ES-MC) & 42.28 \\
        \bottomrule
    \end{tabular}
    \caption{Top 5 y extremos — Caso 1 (2024)}
    \label{tab:caso1_2024}
\end{table}

\subsection{Caso 2: Min-Max + PCA + Media Geométrica}

\begin{table}[H]
    \centering
    \small
    \begin{tabular}{lrr}
        \toprule
        \textbf{Posición} & \textbf{CCAA} & \textbf{Score} \\
        \midrule
        1 & Navarra (ES-NC) & 42.85 \\
        2 & País Vasco (ES-PV) & 38.60 \\
        3 & Cantabria (ES-CB) & 35.42 \\
        4 & La Rioja (ES-LR) & 34.78 \\
        5 & Galicia (ES-GA) & 33.12 \\
        \midrule
        10 & Baleares (ES-IB) & 22.64 \\
        19 & Murcia (ES-MC) & 12.45 \\
        \bottomrule
    \end{tabular}
    \caption{Top 5 y extremos — Caso 2 (2024)}
    \label{tab:caso2_2024}
\end{table}

\subsection{Caso 3: Z-Score + PCA + Media Aritmética}

Para el año 2024, la comparación with Caso 1 (Min-Max + Aritmética) muestra **equivalencia perfecta en todas las posiciones**: $\Delta \text{Rank} = 0$ para las 19 CCAA. Este resultado confirma que under ponderación PCA y agregación lineal, el método de normalización (min-max vs. z-score) no altera el orden relativo de las comunidades autónomas.

Por consiguiente, los rankings del Caso 3 son idénticos a los de Caso 1 para 2024 y para todos los años del período analizado.

\subsection{Caso 4: Min-Max + BoD}

\begin{table}[H]
    \centering
    \small
    \begin{tabular}{lrccc}
        \toprule
        \textbf{Pos.} & \textbf{CCAA} & \textbf{ENV} & \textbf{ECO} & \textbf{Score} \\
        \midrule
        1 & Navarra & 56.4 & 53.0 & 61.4 \\
        2 & Asturias & 67.6 & 33.4 & 59.1 \\
        3 & Galicia & 67.8 & 32.6 & 57.5 \\
        4 & Castilla-León & 63.2 & 37.1 & 56.6 \\
        5 & Extremadura & 62.7 & 39.0 & 56.0 \\
        \bottomrule
    \end{tabular}
    \caption{Top 5 — Caso 4: BoD (2024). Nota: dimensión SOC omitida por espacio}
    \label{tab:caso4_2024}
\end{table}

El método BoD produce un Top 5 identical al Caso 1 en composición pero reordenado internamente. Los scores están normalizados a escala distinta (máximo = 100 en problemas BoD), por lo que los valores numéricos del Caso 4 no son directamente comparables con los anteriores. Lo relevante es la posición relativa.

\section{Caso 4: Min-Max + BoD}

\begin{table}[H]
    \centering
    \small
    \begin{tabular}{lrccc}
        \toprule
        \textbf{Pos.} & \textbf{CCAA} & \textbf{ENV} & \textbf{ECO} & \textbf{Score} \\
        \midrule
        1 & Navarra & 56.4 & 53.0 & 61.4 \\
        2 & Asturias & 67.6 & 33.4 & 59.1 \\
        3 & Galicia & 67.8 & 32.6 & 57.5 \\
        4 & Castilla-León & 63.2 & 37.1 & 56.6 \\
        5 & Extremadura & 62.7 & 39.0 & 56.0 \\
        \bottomrule
    \end{tabular}
    \caption{Top 5 — Caso 4: BoD (2024). Nota: dimensión SOC omitida por espacio}
    \label{tab:caso4_2024}
\end{table}

El método BoD produce un Top 5 identical al Caso 1 en composición pero reordenado internamente. Los scores están normalizados a escala distinta (máximo = 100 en problemas BoD), por lo que los valores numéricos del Caso 4 no son directamente comparables con los anteriores. Lo relevante es la posición relativa.

\section{Caso 5: PCA/FA Rotado (Método OCDE) + Media Aritmética}

El Caso 5 aplica la metodología de \emph{Composite Indicators} del manual OCDE mediante rotación Varimax. Para cada dimensión, se retienen todos los factores con Eigenvalue $> 1$ y se aplica rotación Varimax para producir loadings ortogonales concentrados.

\subsection{Estructura de factores rotados}

\begin{table}[H]
    \centering
    \small
    \begin{tabular}{lc}
        \toprule
        \textbf{Dimensión} & \textbf{Número de factores (Eigenvalue > 1)} \\
        \midrule
        ECO & 7 \\
        SOC & 4 \\
        ENV & 2 \\
        \bottomrule
    \end{tabular}
    \caption{Estructura factorial rotada por dimensión (Caso 5)}
    \label{tab:factores_rotados}
\end{table}

La dimensión económica se desglosa en 7 factores independientes (no uno como en PCA simple PC1):

\begin{enumerate}
    \item Factor ECO-1: Tarificación (cargas altas en precios hoteleros 3--4 estrellas)
    \item Factor ECO-2: Rentabilidad hotelera (RevPAR, ADR)
    \item Factor ECO-3: Ocupación y estancia
    \item Factor ECO-4: Gasto medio del turista
    \item Factor ECO-5: Índices de precios (IPC, IPH)
    \item Factor ECO-6: Empleo turístico
    \item Factor ECO-7: rentabilidad residual / variaciones anuales
\end{enumerate}

Esta descomposición revela que la ``sostenibilidad económica'' del turismo regional no es un constructo unitario sino un agregado de siete dinámicas semindependientes. Una región fuerte en tarificación puede ser débil en ocupación, o viceversa.

\subsection{Pesos factoriales}

En el Caso 5, los pesos de cada indicador se calculan como:

\begin{equation}
    w_j^{\text{OCDE}} = (\text{loading}_j^2 / \sum_{i \in \text{factor}} \text{loading}_i^2) \times (\text{varianza explicada del factor})
\end{equation}

Esto asegura que:
\begin{itemize}
    \item Indicadores altamente correlacionados dentro de un factor comparten peso equitativamente
    \item Factores que explican más varianza contribuyen más al índice
    \item No hay dominación artificial por variables redundantes
\end{itemize}

\subsection{Rankings Caso 5 (2024)}

\begin{table}[H]
    \centering
    \small
    \begin{tabular}{lrr}
        \toprule
        \textbf{Posición} & \textbf{CCAA} & \textbf{Score} \\
        \midrule
        1 & Galicia (ES-GA) & 58.2 \\
        2 & Asturias (ES-AS) & 56.9 \\
        3 & Navarra (ES-NC) & 56.1 \\
        4 & Cantabria (ES-CB) & 55.3 \\
        5 & La Rioja (ES-LR) & 53.8 \\
        \midrule
        15 & Baleares (ES-IB) & 39.4 \\
        19 & Murcia (ES-MC) & 28.6 \\
        \bottomrule
    \end{tabular}
    \caption{Top 5 y extremos — Caso 5 (2024)}
    \label{tab:caso5_2024}
\end{table}

**Diferencias notables respecto a Casos 1--3**:

\begin{itemize}
    \item Galicia lidera (posición 1) frente a Navarra (posición 3) en Caso 1
    \item Asturias sube a posición 2 (versus 2 en Caso 1, equivalente)
    \item Baleares desciende a 15 (equivalente a Caso 1 posición 11, pero score mucho más bajo)
    \item La reordenación es menos drástica que en Caso 2, pero los scores absolutos revelan mayor dispersión entre CCAA
\end{itemize}

La elevación de Galicia en Caso 5 refleja que, bajo factor analysis con rotación, su equilibrio entre los 7 factores económicos es óptimo: no depende de un único super-indicador (como ADR en otras regiones) sino de desempeño robusto multifactorial.

\section{Análisis de volatilidad metodológica

Se calcula la desviación estándar de posiciones entre los cuatro casos para cada CCAA:

\begin{table}[H]
    \centering
    \small
    \begin{tabular}{lc}
        \toprule
        \textbf{CCAA} & \textbf{$\sigma$ de posiciones} \\
        \midrule
        Navarra & 0.0 \\
        País Vasco & 0.5 \\
        La Rioja & 1.2 \\
        Galicia & 0.8 \\
        Baleares & 5.2 \\
        Madrid & 3.1 \\
        Murcia & 0.2 \\
        \bottomrule
    \end{tabular}
    \caption{Volatilidad de posiciones según el método (2024)}
    \label{tab:volatilidad}
\end{table}

Las CCAA del norte (Navarra, País Vasco) exhiben volatilidad **nula o muy baja**, indicativo de liderazgo robusto. En cambio, Baleares presenta $\sigma = 5.2$, lo que refleja sensibilidad extrema: 
- Caso 1: posición 11 (score 46.94)
- Caso 2: posición 10 (score 22.64)
- Caso 4: posición 15 (score no comparable)

Esta volatilidad refleja que Baleares tiene un perfil desequilibrado: muy fuerte en dimensión económica (tarificación hotelera elevada) pero débil en ENV y SOC, de modo que métodos que penalizan desequilibrios (Caso 2) o premian balanceLa (Caso 4) la castigan.

\section{Evolución temporal: trayectorias 2016--2024 (Caso 1)}

La trayectoria de líderes en el período completo bajo Caso 1 (media aritmética):

\begin{table}[H]
    \centering
    \small
    \begin{tabular}{cccccccccc}
        \toprule
        \textbf{Año} & \textbf{2016} & \textbf{2017} & \textbf{2018} & \textbf{2019} & \textbf{2020} & \textbf{2021} & \textbf{2022} & \textbf{2023} & \textbf{2024} \\
        \midrule
        1º & PV & CL & CB & NC & EX & AR & NC & CB & NC \\
        2º & MC & EX & AS & CE & CB & ML & ML & CE & AS \\
        3º & EX & CB & GA & CB & VC & NC & AS & AS & GA \\
        \bottomrule
    \end{tabular}
    \caption{Evolución del Top 3 por año (Caso 1). Códigos: PV=País Vasco, CL=Castilla-León, etc.}
    \label{tab:evol_top3}
\end{table}

**Observaciones**:
\begin{itemize}
    \item \textbf{Navarra (ES-NC)}: lidera en 2024 y 2022, con apariciones fuertes en múltiples años. Consistencia elevada.
    \item \textbf{Castilla-León (ES-CL)}: fuerte en 2017, luego desciende. Posición 4 en 2024.
    \item \textbf{Cantabria (ES-CB)}: volatilidad: liderato en 2018 y 2023, pero posición media en años intermedios.
    \item \textbf{Murcia (ES-MC)}: consistentemente en posiciones bajas desde 2016 (posición 19 en 2024).
\end{itemize}

El orden de cambio estructuralmente no es aleatorio: regiones del norte y noroeste (Galicia, Asturias, Cantabria) muestran trayectoria ascendente, mientras que Murcia exhibe estancamiento en los últimos puestos.

\section{Aplicación del método BoD a nivel dimensional}

El método BoD aplicado sobre los tres scores dimensionales (en lugar de los 52 indicadores) revela **pesos óptimos heterogéneos por CCAA**:

\begin{table}[H]
    \centering
    \small
    \begin{tabular}{lccc}
        \toprule
        \textbf{CCAA} & \textbf{$w_{\text{ENV}}$} & \textbf{$w_{\text{ECO}}$} & \textbf{$w_{\text{SOC}}$} \\
        \midrule
        Navarra & 0.30 & 0.25 & 0.45 \\
        Baleares & 0.05 & 0.90 & 0.05 \\
        Murcia & 0.15 & 0.10 & 0.75 \\
        \bottomrule
    \end{tabular}
    \caption{Pesos BoD elegidos óptimamente por CCAA (ejemplo)}
    \label{tab:pesos_bod_ccaa}
\end{table}

**Interpretación**:
\begin{itemize}
    \item Navarra elige pesos equilibrados con énfasis en SOC (45\%), reflejando su fortaleza social
    \item Baleares se especializa en ECO (90\%), minimizando el peso de dimensiones donde es débil
    \item Murcia enfatiza SOC (75\%), tratando de compensar debilidades en ECO y ENV
\end{itemize}

Estos pesos óptimos demuestran que bajo BoD, el ranking recompensa la capacidad de especialización, no el balance multidimensional.

\section{Síntesis comparativa de casos (2024)}

\begin{table}[H]
    \centering
    \footnotesize
    \begin{tabular}{lcccc}
        \toprule
        \textbf{CCAA} & \textbf{Caso 1} & \textbf{Caso 2} & \textbf{Caso 3} & \textbf{Caso 4} \\
        \midrule
        Navarra & 1 & 1 & 1 & 1 \\
        País Vasco & 6 & 2 & 6 & 7 \\
        La Rioja & 9 & 4 & 9 & 10 \\
        Baleares & 11 & 10 & 11 & 15 \\
        Madrid & 14 & 8 & 14 & 12 \\
        \bottomrule
    \end{tabular}
    \caption{Comparativa de rankings por caso (CCAA seleccionadas, 2024)}
    \label{tab:resumen_casos}
\end{table}

**Patrones observados**:
\begin{itemize}
    \item Navarra domina en todos los casos (posición 1 invariante)
    \item País Vasco gana posiciones en Caso 2 (media geométrica)
    \item Baleares pierde posiciones en Casos 2 y 4 (penalización de desequilibrios)
    \item Casos 1 y 3 dan resultados idénticos (normalización irrelevante)
\end{itemize}

\end{document}
